\section*{الفصل \ref{ch:g11:polarity-and-cohesion} --- الكهرسلبية والقطبية والقوى بين الجزيئات}

\begin{solution}{exo:g11:polarity-and-cohesion:1}
\ce{H-Cl}: الكلور (3.16 مقابل 2.20). \ce{O-H}: الأكسجين (3.44).
\ce{C-O}: الأكسجين (3.44 مقابل 2.55).
\end{solution}

\begin{solution}{exo:g11:polarity-and-cohesion:2}
\ce{H-H}: 0، غير قطبية. \ce{C-H}: 0.35، تُعدّ غير قطبية. \ce{O-H}:
1.24، قطبية. \ce{N-H}: 0.84، قطبية. \ce{Cl-Cl}: 0، غير قطبية. \ce{C-Cl}:
0.61، قطبية.
\end{solution}

\begin{solution}{exo:g11:polarity-and-cohesion:3}
النيتروجين (3.04 > 2.19) والأكسجين (3.44 > 2.58): تتناقص \omterm{def:g11:polarity-and-cohesion:electronegativity}{الكهرسلبية} نزولًا في
\omterm{def:g9:periodic-table-first-look:period}{الزمرة}. والنيتروجين (3.04 > 2.55) والمغنيسيوم
(1.31 > 0.93): تتزايد من اليسار إلى اليمين على امتداد \omterm{def:g9:periodic-table-first-look:period}{الدورة}.
\end{solution}

\begin{solution}{exo:g11:polarity-and-cohesion:4}
الأمونياك والماء: لكل منهما \omterm{def:g7:atoms-and-molecules:atom}{ذرات} هيدروجين مرتبطة بالنيتروجين أو الأكسجين، وأزواج
حرة على تلك \omterm{def:g7:atoms-and-molecules:atom}{الذرة}. أما الميثان فليس فيه هيدروجين على N أو O أو F؛ وكلوريد
الهيدروجين فيه هيدروجين مرتبط بالكلور، وهو ليس من \omterm{def:g7:atoms-and-molecules:atom}{الذرات} الثلاث الواردة في
التعريف.
\end{solution}

\begin{solution}{exo:g11:polarity-and-cohesion:5}
نعم: فتآثرات فان دير فالس موجودة بين كل \omterm{def:g7:atoms-and-molecules:molecule}{الجزيئات}. لكن \omterm{def:g7:atoms-and-molecules:molecule}{جزيئات} الميثان صغيرة (10
\omterm{def:g9:inside-the-atom:nucleus}{إلكترونات})، فتكون هذه التآثرات ضعيفة جدًا وتنكسر عند درجة حرارة منخفضة جدًا: فالميثان
يغلي قرب
\qty{-162}{\celsius}.
\end{solution}

\begin{solution}{exo:g11:polarity-and-cohesion:6}
\ce{CH2Cl2}: قطبي (رابطتان \ce{C-Cl} قطبيتان ورابطتان \ce{C-H} شبه غير
قطبيتين: ولا تناظر يلغيها). \ce{CCl4}: غير قطبي (أربع
روابط متماثلة على رؤوس رباعي أوجه). \ce{NH3}: قطبي
(هرم). \ce{CO2}: غير قطبي (خطي، رابطتان متعاكستان). \ce{BF3}:
غير قطبي (ثلاث روابط قطبية متماثلة على \ang{120} في مستوٍ تتلاشى).
\end{solution}

\begin{solution}{exo:g11:polarity-and-cohesion:7}
يمكن أن يُعطى هيدروجين المجموعة \ce{O-H} (لا هيدروجينات
\ce{CH3}، المرتبطة بالكربون)؛ وتستطيع \omterm{def:g7:atoms-and-molecules:atom}{ذرة} الأكسجين، بزوجيها الحرين، أن
تستقبل. نعم: فالمجموعة \ce{O-H} في الميثانول تستطيع أن ترتبط بأكسجين \omterm{def:g7:atoms-and-molecules:molecule}{جزيء}
ماء، والمجموعة \ce{O-H} في الماء بأكسجين الميثانول، مثلًا
\ce{CH3-O-H}\,$\cdots$\,\ce{OH2}.
\end{solution}

\begin{solution}{exo:g11:polarity-and-cohesion:8}
\omterm{def:g7:atoms-and-molecules:molecule}{الجزيئات} الثلاثة غير قطبية ومن الشكل نفسه، لكنها أكبر
فأكبر (10 و18 و36 إلكترونًا): فتكبر تآثرات فان دير فالس بينها،
وترتفع معها درجات غليانها.
\end{solution}

\begin{solution}{exo:g11:polarity-and-cohesion:9}
\omterm{def:g9:periodic-table-first-look:period}{الزمرة} 14: فالميثان ليس فيه هيدروجين مرتبط \omterm{def:g7:atoms-and-molecules:atom}{بذرة} N أو O أو F ولا يكوّن روابط
هيدروجينية، فهو يتبع اتجاه زمرته.
\end{solution}

\begin{solution}{exo:g11:polarity-and-cohesion:10}
تآثرات فان دير فالس: \omterm{def:g7:atoms-and-molecules:molecule}{فالجزيئات} تكبر من \ce{HCl} إلى
\ce{HI}، وازدياد هذه التآثرات مع الحجم يفوق نقص التجاذب بين الشحنات
الجزئية.
\end{solution}

\begin{solution}{exo:g11:polarity-and-cohesion:11}
\omterm{def:g7:atoms-and-molecules:molecule}{جزيء} اليود أكبر بكثير من \omterm{def:g7:atoms-and-molecules:molecule}{جزيء} الكلور (106
\omterm{def:g9:inside-the-atom:nucleus}{إلكترونات} مقابل 34): فتآثرات فان دير فالس بين جزيئاته أقوى بكثير، وتكفي
لإمساك \omterm{def:g7:atoms-and-molecules:molecule}{الجزيئات} في جسم صلب في درجة حرارة الغرفة.
\end{solution}

\begin{solution}{exo:g11:polarity-and-cohesion:12}
البروبان غير قطبي: تآثرات فان دير فالس فقط. والميثوكسي ميثان
قطبي (رابطتان \ce{C-O} قطبيتان في \ce{C-O-C} منحنية) لكن ليس على أكسجينه
هيدروجين: فلا روابط هيدروجينية بين جزيئاته. والإيثانول
قطبي ومجموعاته \ce{O-H} تكوّن روابط هيدروجينية. وكلما اشتد
التجاذب ارتفعت درجة الغليان: البروبان < الميثوكسي ميثان
< الإيثانول.
\end{solution}

\begin{solution}{exo:g11:polarity-and-cohesion:13}
الشبكة المفتوحة للجليد تشغل حيزًا أكبر من \omterm{def:g7:atoms-and-molecules:molecule}{الجزيئات} نفسها في السائل، فالجليد
أقل كثافة من الماء السائل ويطفو. والبحيرة تتجمد من أعلاها: فتعزل طبقة الجليد
الماء الذي تحتها، فيبقى سائلًا، وتنجو الأسماك.
\end{solution}

\begin{solution}{exo:g11:polarity-and-cohesion:14}
الميثان الصلب (تآثرات فان دير فالس بين \omterm{def:g7:atoms-and-molecules:molecule}{جزيئات} صغيرة) <
الجليد (روابط هيدروجينية) < كلوريد البوتاسيوم (تجاذب بين \omterm{def:g9:ions:ion}{الأيونات}).
\end{solution}

\begin{solution}{exo:g11:polarity-and-cohesion:15}
\omterm{def:g7:atoms-and-molecules:molecule}{لجزيء} الماء ذرتا هيدروجين يعطيهما وزوجان حران يستقبل بهما: فيستطيع كل
\omterm{def:g7:atoms-and-molecules:molecule}{جزيء} أن يشارك في أربع روابط هيدروجينية، اثنتان معطاة واثنتان مستقبلة، وهذا يبني
شبكة كاملة. أما \omterm{def:g7:atoms-and-molecules:molecule}{جزيء} فلوريد الهيدروجين فله ثلاثة أزواج حرة لكن \omterm{def:g7:atoms-and-molecules:atom}{ذرة} هيدروجين واحدة
فقط: فهو يعطي \omterm{def:g11:polarity-and-cohesion:hydrogen-bond}{رابطة هيدروجينية} واحدة، فلا تتكوّن في المتوسط إلا رابطة واحدة لكل
\omterm{def:g7:atoms-and-molecules:molecule}{جزيء} (سلاسل متعرجة). فللماء نحو ضعف الروابط الهيدروجينية لكل \omterm{def:g7:atoms-and-molecules:molecule}{جزيء}،
فيغلي عند درجة أعلى.
\end{solution}

\begin{solution}{pb:g11:polarity-and-cohesion:1}
\textbf{1.} \ce{H-O-H} و\ce{H-S-H} و\ce{H-Se-H}، ولكل \omterm{def:g7:atoms-and-molecules:atom}{ذرة} مركزية
زوجان حران: فلكل من O وS وSe ستة \omterm{def:g9:inside-the-atom:nucleus}{إلكترونات} تكافؤ، لأنها في
\omterm{def:g9:periodic-table-first-look:period}{الزمرة} نفسها.

\textbf{2.} $M(\ce{H2O}) = \qty{18.0}{g/mol}$، $M(\ce{H2S}) =
\qty{34.1}{g/mol}$، $M(\ce{H2Se}) = \qty{81.0}{g/mol}$.

\textbf{3.} $3.44 - 2.20 = 1.24$؛ $2.58 - 2.20 = 0.38$؛
$2.55 - 2.20 = 0.35$. فالرابطة \ce{O-H} وحدها قطبية بوضوح.

\textbf{4.} منحنية (أربع مجموعات، منها زوجان حران). الماء قطبي؛ أما كبريتيد
الهيدروجين وسيلينيد الهيدروجين فقطبيتهما ضئيلة جدًا.

\textbf{5.} 10 و18 و36 إلكترونًا: فسيلينيد الهيدروجين له أقوى
تآثرات فان دير فالس.

\textbf{6.} الماء وحده: \omterm{def:g7:atoms-and-molecules:atom}{فذرات} الهيدروجين فيه مرتبطة بالأكسجين، أحد
\omterm{def:g7:atoms-and-molecules:atom}{الذرات} الثلاث الشديدة \omterm{def:g11:polarity-and-cohesion:electronegativity}{الكهرسلبية}؛ أما الكبريت والسيلينيوم فليسا من \omterm{def:g11:polarity-and-cohesion:electronegativity}{الكهرسلبية}
بحيث يعطيان \omterm{def:g7:atoms-and-molecules:atom}{ذرات} الهيدروجين المرتبطة بهما $\delta^+$
واضحة.

\textbf{7.} أربع: اثنتان عن طريق ذرتي الهيدروجين فيه، واثنتان يستقبلهما
زوجاه الحران.

\textbf{8.} تآثرات فان دير فالس (بما فيها تجاذب ضعيف
بين الشحنات الجزئية الصغيرة).

\textbf{9.} $212.9 - 273.15 = \qty{-60.3}{\celsius}$
و\qty{-41.3}{\celsius}.

\textbf{10.} $-41.3 - (-60.3) = \qty{19.0}{\celsius}$.

\textbf{11.} \omterm{def:g7:atoms-and-molecules:molecule}{الجزيء} أكبر، فتكون تآثرات فان دير فالس بين جزيئاته
أقوى.

\textbf{12.} $-60.3 - 19.0 = \qty{-79.3}{\celsius}$.

\textbf{13.} الارتفاع $-88.1 - (-112) = \qty{23.9}{\celsius}$؛ والتنبؤ
للميثان $-112 - 23.9 = \qty{-136}{\celsius}$، مقابل \qty{-162}{\celsius}
في الواقع: فالاستقراء أعلى بنحو \qty{26}{\celsius}. فهو غير
دقيق؛ والاتجاه الحقيقي ليس خطًا مستقيمًا.

\textbf{14.} الارتفاع $\qty{18.7}{\celsius}$؛ والتنبؤ لفلوريد الهيدروجين
$-85.1 - 18.7 = \qty{-103.8}{\celsius}$، مقابل \qty{19.6}{\celsius}: أي
فجوة نحو \qty{123}{\celsius}.

\textbf{15.} $100 - (-79.3) = \qty{179}{\celsius}$، أي نحو
\qty{180}{\celsius}.

\textbf{16.} الارتفاع $\qty{25.3}{\celsius}$؛ والتنبؤ $-87.8 - 25.3 =
\qty{-113}{\celsius}$؛ والفجوة $-33.4 - (-113) = \qty{80}{\celsius}$.

\textbf{17.} الماء (\qty{180}{\celsius}) > فلوريد الهيدروجين
(\qty{123}{\celsius}) > الأمونياك (\qty{80}{\celsius}). فالماء يعطي ذرتي
هيدروجين ويستقبل بزوجين حرين: أربع روابط هيدروجينية لكل
\omterm{def:g7:atoms-and-molecules:molecule}{جزيء}، مستعملة كلها. ولفلوريد الهيدروجين هيدروجين واحد يعطيه، وللأمونياك
\omterm{def:g10:lewis-and-shape:lone-pair}{زوج حر} واحد فقط يستقبل به: ففي المتوسط روابط أقل لكل
\omterm{def:g7:atoms-and-molecules:molecule}{جزيء}، والنيتروجين، وهو أقل \omterm{def:g11:polarity-and-cohesion:electronegativity}{كهرسلبية}، يكوّن روابط أضعف.

\textbf{18.} غازًا: فكان سيغلي قرب \qty{-80}{\celsius}. ولما كان على الأرض
ماء سائل، ولا محيطات، ولا مطر، ولا شيء من الحياة
التي تعتمد عليها.

\textbf{19.} فهو يستقرئ خطًا مستقيمًا مرسومًا عبر نقطتين
فقط، والاختبار على \omterm{def:g9:periodic-table-first-look:period}{الزمرة} 14 يبيّن أن مثل هذا الاستقراء قد يخطئ
بنحو \qty{25}{\celsius}.

\textbf{20.} لولا الروابط الهيدروجينية لغلى الماء عند نحو
\qty{-80}{\celsius}، أي أدنى من درجة غليانه الحقيقية بنحو
\qty{180}{\celsius}.
\end{solution}
